%TEX root=../main.tex
\section{Related Work}

\paragraph{Calibrated probability estimates} are considered necessary for empirical risk
analysis tools \cite{berk2017fairness,crowson-calibration-risk-scores,dieterich-northpointe-fairness,flores-re-propublica-fair}.
In practical applications, uncalibrated probability estimates can be misleading
in the sense that the end user of these estimates has an incentive to mistrust (and therefore potentially misuse) them.
We note however that calibration does not remove all potential for misuse, as the end user's biases
might cause her or him to treat estimates differently based on group membership.
%
There are several post-processing methods for producing calibrated outputs from classification
algorithms. For example, Platt Scaling \cite{platt1999probabilistic} passes outputs through
a learned sigmoid function, transforming them into calibrated probabilities.
Histogram Binning and Isotonic Regression \cite{zadrozny2001obtaining}
learn a general monotonic function from outputs to probabilities.
See \cite{niculescu2005predicting} and \cite{guo2017calibration} for empirical
comparisons
of these methods.

\paragraph{Equalized Odds}~\cite{hardt2016equality},
also referred to as \emph{Disparate Mistreatment} \cite{zafar2017fairness},
ensures that no error type disproportionately affects any particular group.
\citet{hardt2016equality} provide a post-processing technique to achieve this
framework, while \citet{zafar2017fairness} introduce optimization constraints
to achieve non-discrimination at training time.
Recently, this framework has received significant attention from the algorithmic
fairness community.
Researchers have found that it is incompatible with other notions
of fairness \cite{kleinberg2016inherent,chouldechova2017fair,corbett2017algorithmic}.
Additionally, \citet{woodworth2017learning} demonstrate that, under certain assumptions, post-processing
methods for achieving non-discrimination may be suboptimal.

\paragraph{Alternative fairness frameworks} exist and are continuously proposed.
%As the field of algorithm fairness evolves rapidly, there are
%many alternative frameworks of machine learning fairness.
We highlight several of these works, though by no means offer
a comprehensive list. (More thorough reviews can be found in
\cite{barocas2016big,berk2017fairness,romei2014multidisciplinary}).
It has been shown that, under most frameworks of fairness, there is a trade-off
between algorithmic performance and non-discrimination
\cite{berk2017fairness,corbett2017algorithmic,hardt2016equality,zliobaite2015relation}.
Several works approach fairness through the lens of
\emph{Statistical Parity}
\cite{calders2009building,kamiran2009classifying,calders2010three,kamishima2011fairness,zemel2013learning,louizos2015variational,edwards2015censoring,johndrow2017algorithm}.
Under this definition, group membership should not affect the prediction of a classifier,
i.e. members of different groups should have the same probability of receiving
a positive-class prediction.
%Statistical parity can be achieved through
%regularization mechanisms \cite{kamishima2011fairness}, optimization
%constraints \cite{zafar2015learning}, or by feature representations
%that remove all group-specific information
%\cite{zemel2013learning,louizos2015variational,edwards2015censoring}.
%
However,
it has been argued that Statistical Parity may not be applicable
in many scenarios
\cite{chouldechova2017fair,dwork2012fairness,hardt2016equality,kleinberg2016inherent},
as it
%Statistical parity essentially amounts to affirmative action,
%and has the potential to significantly harm accuracy in scenarios where group
%membership is strongly correlated with outcome. This is especially harmful in
%medical diagnosis and risk-assessment scenarios where errors are costly and to be avoided.
%
attempts to guarantee equal representation. For example, it is
inappropriate in criminal justice, where base rates differ
across different groups.
%
A related notion is \emph{Disparate Impact}
\cite{feldman2015certifying,zafar2015learning}, which states that the
prediction rates for any two groups should not differ by more than $80\%$
(a number motivated by legal requirements).
%
\citet{dwork2012fairness} introduce a notion of fairness based on
the idea that similar individuals should receive similar outcomes,
though it challenging to achieve this notion in practice.
%Though this is
%a rich definition of fairness, it is challenging to implement in practice
%because it requires an externally defined similarity metric between individuals.
Fairness has also been considered in online learning
\cite{joseph2016fairness,kearns2017meritocratic}, unsupervised learning \cite{bolukbasi2016man},
and causal inference \cite{kilbertus2017avoiding,kusner2017counterfactual}.

